\section{Methodology}
\label{sec:method}

\subsection{Pipeline Overview}
SmartPlate processes a single RGB plate photograph through four sequential stages:
(1) semantic segmentation into 103 fine-grained food classes,
(2) aggregation of per-pixel class predictions into 7 Harvard food-group proportions,
(3) health scoring via a closed-form rubric,
and (4) macro-nutrient estimation via a per-group USDA lookup table.
The segmentation stage is the only learned component; stages 2--4 are deterministic functions of the pixel-area proportions, making the final score fully auditable.

\subsection{Dataset}
We train and evaluate on \textbf{FoodSeg103}~\cite{wu2021foodseg103}, a publicly available benchmark of $7{,}118$ food images with dense pixel-level masks across $103$ ingredient categories (plus a background class), split into $4{,}983$ training and $2{,}135$ validation images.
Images are variable in resolution; we resize all inputs to $512\times512$ during training.
FoodSeg103 covers a wide variety of cuisines, making it more appropriate than single-cuisine datasets for the mixed-plate setting SmartPlate targets.

\subsection{Segmentation Model}
We fine-tune \textbf{SegFormer-B0}~\cite{xie2021segformer}, initialized from the ADE20K-pretrained checkpoint \texttt{nvidia/segformer-b0-finetuned-ade-512-512}.
SegFormer couples a hierarchical Mix Transformer (MiT-B0) encoder with a lightweight all-MLP decoder head.
The ADE20K classification head (150-way) is replaced with a 104-way head (103 food classes + background), and all parameters are updated during fine-tuning.
The resulting model has $3.74$ M trainable parameters.

\paragraph{Training details.}
We train for $5$ epochs with the AdamW optimizer ($\text{lr}=6\times10^{-5}$, weight decay $= 10^{-4}$, batch size $= 2$).
Training augmentation is a random horizontal flip applied jointly to the image and mask.
We use the SegFormer's built-in cross-entropy loss computed over the upsampled logit map.
Training ran on a Google Colab T4 GPU in approximately $40$ minutes; inference runs on Apple Silicon (PyTorch MPS backend).
The best checkpoint (by 103-class validation mIoU) is saved and used for all downstream evaluation.

\subsection{Class-to-Group Aggregation}
\label{sec:mapping}
The $103$ FoodSeg103 classes are manually mapped to $7$ food groups drawn from the Harvard Healthy Eating Plate~\cite{harvardplate}:
\textit{vegetables}, \textit{fruits}, \textit{whole grains}, \textit{healthy protein}, \textit{refined grains}, \textit{red/processed meat}, and \textit{sugary/fatty foods}.
Eight classes (beverages, sauce, soup, and \textit{other ingredients}) do not correspond to a solid food group and are assigned to background, excluded from the denominator.
In total, $96$ of $104$ classes (including background) are mapped.
Given a predicted $H \times W$ class-id mask, the proportion of food area belonging to group $g$ is:
\begin{equation}
    p_g = \frac{\sum_{(i,j)} \mathbf{1}[\hat{y}_{ij} \in \mathcal{C}_g]}{\sum_{(i,j)} \mathbf{1}[\hat{y}_{ij} \notin \mathcal{C}_\text{bg}]},
\end{equation}
where $\mathcal{C}_g$ is the set of class ids assigned to group $g$ and $\mathcal{C}_\text{bg}$ is the background set.
The proportions $\{p_g\}$ sum to $1$ over the $7$ food groups.

\subsection{Health Scoring}
\label{sec:scoring}
The health score $s \in [0, 100]$ is computed from the food-group proportions via a three-term formula:
\begin{align}
    \text{base} &= 100 \cdot \left(\sum_{g \in \mathcal{G}^+} p_g + \alpha \cdot p_\text{red meat}\right), \label{eq:base}\\
    \text{balance penalty} &= \text{base} \cdot 0.15 \cdot \frac{1}{2}\sum_{g \in \mathcal{G}^+}|p_g - t_g|, \label{eq:balance}\\
    \text{negative penalty} &= 100 \cdot \sum_{g \in \mathcal{G}^-} \lambda_g \cdot p_g, \label{eq:neg}\\
    s &= \operatorname{clip}(\text{base} - \text{balance penalty} - \text{negative penalty},\; 0,\; 100). \label{eq:score}
\end{align}
Here $\mathcal{G}^+ = \{\text{vegetables, fruits, whole grains, healthy protein}\}$ are the positive groups with Harvard target proportions $t_g \in \{0.30, 0.20, 0.25, 0.25\}$;
$\mathcal{G}^- = \{\text{refined grains, red/processed meat, sugary/fatty}\}$ are the negative groups;
$\alpha = 0.50$ gives red meat partial protein credit; and
the per-group penalty strengths are $\lambda_\text{refined grains} = 0.15$, $\lambda_\text{red meat} = 0.35$, $\lambda_\text{sugary/fatty} = 0.80$.
A plate matching the Harvard target proportions exactly with no negative groups scores $100$; an all-junk plate scores near $0$.
The balance penalty (Eq.~\ref{eq:balance}) uses a normalized $L_1$ distance to the ideal proportions, reducing the score by up to $15\%$ of the base when the positive groups are maximally imbalanced.

\subsection{Nutritional Estimation}
Given the food-group proportions $\{p_g\}$ and an assumed total plate mass of $400\,$g, we estimate aggregate macros as:
\begin{equation}
    M_k = \sum_g p_g \cdot 400 \cdot \frac{m_{g,k}}{100},
\end{equation}
where $m_{g,k}$ is the per-100\,g macro value (kcal, protein, carbohydrate, or fat) for group $g$ from USDA FoodData Central~\cite{usdafdc}.
The $400\,$g assumption is a fixed prior because a monocular RGB image contains no depth information; we report macros as coarse estimates and discuss this limitation in Section~\ref{sec:conclusion}.

\subsection{ResNet-50 Multi-Label Baseline}
To quantify what pixel-level segmentation buys over image-level classification, we train a \textbf{ResNet-50}~\cite{he2016resnet} multi-label classifier.
The ImageNet-pretrained backbone's fully-connected head is replaced with a 7-way linear layer followed by sigmoid activations, one output per Harvard food group.
A group is labeled as present (positive) in a training image if its ground-truth pixel proportion exceeds $5\%$ of the food area; the prediction threshold at inference is $0.40$.
The model is trained for $10$ epochs with AdamW ($\text{lr}=10^{-4}$, weight decay $= 10^{-4}$) and a cosine annealing scheduler, using binary cross-entropy loss.
Crucially, this model predicts \emph{which} food groups are present but not \emph{how much} of each, so it cannot produce a health score --- establishing the necessity of pixel-level proportions for this task.
